\documentclass[11pt,a4paper]{article}
\usepackage[T1]{fontenc}
\usepackage[utf8]{inputenc}
\usepackage{amsmath,amssymb,xcolor}
\usepackage[margin=25mm]{geometry}
\usepackage[colorlinks=true,urlcolor=blue]{hyperref}
\definecolor{nexorablue}{HTML}{2459B5}
\setlength{\parindent}{0pt}
\setlength{\parskip}{7pt}
\setlength{\emergencystretch}{2em}
\renewcommand{\familydefault}{\sfdefault}
\begin{document}

{\small\bfseries\color{nexorablue} NEXORA CPGE}\par
{\small MP / NEX-IND-05 / Énoncé}\par
{\LARGE\bfseries Courant transitoire dans une bobine secondaire}\par
\vspace{0.25em}\hrule\vspace{0.6em}

Deux bobines sont couplées avec $M=0{,}040\,\mathrm H$ dans les orientations choisies. Le courant primaire croît avec $\dot i_1=5{,}0\,\mathrm{A\,s^{-1}}$ dès $t=0$. La boucle secondaire a $R=3{,}0\,\mathrm\Omega$, $L=0{,}12\,\mathrm H$ et un courant initial nul. Trouver $i_2(t)$ pendant la rampe et l’évaluer à $80\,\mathrm{ms}$.


\end{document}
